% Silabus Mata Kuliah Kriptografi - OBE
% Program Studi Matematika S1
% 3 SKS, 16 pertemuan, implementasi dengan bahasa C

\documentclass[12pt,a4paper]{article}
\usepackage[utf8]{inputenc}
\usepackage[T1]{fontenc}
\usepackage[indonesian]{babel}
\usepackage{geometry}
\usepackage{longtable}
\usepackage{array}
\usepackage{booktabs}
\usepackage{enumitem}
\usepackage[hidelinks,breaklinks]{hyperref}
\geometry{margin=2.5cm}

\begin{document}

\begin{center}
\textbf{\Large RENCANA PEMBELAJARAN SEMESTER (RPS)}\\
\textbf{Berbasis Outcome-Based Education (OBE)}\\[0.5cm]
\textbf{PROGRAM STUDI MATEMATIKA}\\
\textbf{FAKULTAS MIPA / SAINS}\\
\textbf{MATA KULIAH KRIPTOGRAFI}\\
\end{center}

\vspace{1cm}

% ============================================================
% 1. IDENTITAS MATA KULIAH
% ============================================================
\section*{1. Identitas Mata Kuliah}
\begin{tabular}{ll}
Nama Program Studi & : Matematika \\
Nama Mata Kuliah & : Kriptografi \\
Kode Mata Kuliah & : MAT-XXX \\
Semester & : Ganjil/Genap (sesuai kurikulum) \\
SKS / Bobot Kredit & : 3 SKS (2 Teori, 1 Praktikum) \\
Jumlah Pertemuan & : 16 pertemuan \\
Dosen Pengampu & : [Nama Dosen] \\
Prasyarat & : Aljabar Linear, Matematika Diskrit \\
Tanggal Penyusunan & : [Tanggal] \\
\end{tabular}

\vspace{0.5cm}

% ============================================================
% 2. CAPAIAN PEMBELAJARAN LULUSAN (CPL)
% ============================================================
\section*{2. Capaian Pembelajaran Lulusan (CPL)}

CPL yang dibebankan pada mata kuliah ini mencakup kompetensi lulusan dalam aspek pengetahuan, keterampilan, dan sikap:

\begin{itemize}[leftmargin=*]
  \item \textbf{CPL-1 (Pengetahuan):} Menguasai konsep teoretis matematika termasuk aljabar, analisis, dan matematika diskrit untuk menyelesaikan masalah.
  
  \item \textbf{CPL-2 (Pengetahuan):} Menguasai prinsip-prinsip komputasi dan pemrograman untuk implementasi solusi matematika.
  
  \item \textbf{CPL-3 (Keterampilan Umum):} Mampu menerapkan pemikiran logis, kritis, sistematis, dan inovatif dalam konteks pengembangan ilmu pengetahuan.
  
  \item \textbf{CPL-4 (Keterampilan Khusus):} Mampu merancang dan mengimplementasikan algoritma matematika menggunakan bahasa pemrograman.
  
  \item \textbf{CPL-5 (Keterampilan Khusus):} Mampu menganalisis dan menyelesaikan masalah matematika terapan.
  
  \item \textbf{CPL-6 (Sikap):} Menunjukkan sikap bertanggung jawab atas pekerjaan di bidang keahliannya secara mandiri dan mampu bekerja sama dalam tim.
\end{itemize}

\vspace{0.5cm}

% ============================================================
% 3. CAPAIAN PEMBELAJARAN MATA KULIAH (CPMK)
% ============================================================
\section*{3. Capaian Pembelajaran Mata Kuliah (CPMK)}

Kemampuan atau kompetensi spesifik yang diharapkan mahasiswa kuasai setelah menyelesaikan mata kuliah:

\begin{itemize}[leftmargin=*]
  \item \textbf{CPMK-1:} Mahasiswa mampu menjelaskan prinsip-prinsip dasar kriptografi dan konsep keamanan informasi (confidentiality, integrity, authenticity).
  
  \item \textbf{CPMK-2:} Mahasiswa mampu menerapkan aritmetika modular, algoritma Euclid, dan teori bilangan dalam konteks kriptografi.
  
  \item \textbf{CPMK-3:} Mahasiswa mampu menganalisis dan mengimplementasikan algoritma kriptografi klasik (Caesar, Vigenère, substitusi) dalam bahasa C.
  
  \item \textbf{CPMK-4:} Mahasiswa mampu menganalisis dan mengimplementasikan algoritma kriptografi modern simetris (AES, DES) dan asimetris (RSA) dalam bahasa C.
  
  \item \textbf{CPMK-5:} Mahasiswa mampu menerapkan fungsi hash dan tanda tangan digital dalam keamanan data.
  
  \item \textbf{CPMK-6:} Mahasiswa mampu merancang dan menganalisis protokol kriptografi serta manajemen kunci.
  
  \item \textbf{CPMK-7:} Mahasiswa mampu melakukan kriptanalisis dasar terhadap cipher sederhana (frequency analysis, Kasiski examination).
  
  \item \textbf{CPMK-8:} Mahasiswa mampu mengimplementasikan algoritma kriptografi menggunakan bahasa pemrograman C.
\end{itemize}

\vspace{0.5cm}

% ============================================================
% 4. SUB-CPMK / INDIKATOR PENCAPAIAN
% ============================================================
\section*{4. Sub-CPMK / Indikator Pencapaian}

Penjabaran CPMK menjadi indikator yang lebih terukur dan dapat diuji:

\begin{itemize}[leftmargin=*]
  \item \textbf{Sub-CPMK 1.1:} Menjelaskan model komunikasi aman dan ancaman keamanan informasi
  \item \textbf{Sub-CPMK 1.2:} Mendefinisikan confidentiality, integrity, dan authenticity
  \item \textbf{Sub-CPMK 1.3:} Membedakan kriptografi simetris dan asimetris
  \item \textbf{Sub-CPMK 2.1:} Menghitung operasi aritmetika modular (penjumlahan, perkalian, invers) dan mengimplementasikan algoritma Euclid diperluas dalam C
  \item \textbf{Sub-CPMK 2.2:} Menentukan bilangan prima, uji primalitas, dan menghitung pangkat modular (fast exponentiation) dalam C
  \item \textbf{Sub-CPMK 3.1:} Mengimplementasikan cipher Caesar dan Vigenère dalam bahasa C
  \item \textbf{Sub-CPMK 3.2:} Menganalisis kelemahan cipher klasik dan melakukan enkripsi/dekripsi dengan cipher substitusi
  \item \textbf{Sub-CPMK 4.1:} Menjelaskan struktur AES dan DES serta mode operasi block cipher (ECB, CBC)
  \item \textbf{Sub-CPMK 4.2:} Mengimplementasikan algoritma RSA dalam bahasa C (dengan bilangan kecil atau library)
  \item \textbf{Sub-CPMK 5.1:} Menjelaskan sifat fungsi hash kriptografis dan mengimplementasikan SHA/MD5 untuk integritas data
  \item \textbf{Sub-CPMK 5.2:} Merancang tanda tangan digital berbasis RSA
  \item \textbf{Sub-CPMK 6.1:} Menganalisis protokol Diffie-Hellman dan siklus hidup manajemen kunci
  \item \textbf{Sub-CPMK 7.1:} Melakukan frequency analysis dan Kasiski examination pada cipher substitusi
  \item \textbf{Sub-CPMK 8.1:} Membuat dokumentasi dan pengujian kode C untuk implementasi kriptografi
\end{itemize}

\vspace{0.5cm}

% ============================================================
% 5. MATERI PEMBELAJARAN (BAHAN KAJIAN)
% ============================================================
\section*{5. Materi Pembelajaran (Bahan Kajian)}

Daftar topik materi yang relevan dengan Sub-CPMK dan CPMK:

\begin{enumerate}[leftmargin=*]
  \item Pengantar Kriptografi: model keamanan, terminologi, sejarah
  \item Aritmetika modular dan operasi dasar (penjumlahan, perkalian, invers modular)
  \item Algoritma Euclid dan Euclid diperluas
  \item Teori bilangan: bilangan prima, uji primalitas, teorema Fermat dan Euler
  \item Pangkat modular (fast exponentiation)
  \item Algoritma kriptografi klasik: Caesar, substitusi monoalfabetik, Vigenère, transposisi
  \item Block cipher: DES, AES, mode operasi ECB, CBC, CFB, OFB
  \item Kriptografi asimetris: RSA (pembangkitan kunci, enkripsi, dekripsi)
  \item Fungsi hash (SHA, MD5), integritas data
  \item Tanda tangan digital dan sertifikat digital
  \item Protokol kriptografi: Diffie-Hellman, TLS dasar
  \item Manajemen kunci: siklus hidup, distribusi kunci
  \item Kriptanalisis: frequency analysis, Kasiski examination, serangan known-plaintext
\end{enumerate}

\subsection*{Rencana Pembelajaran 16 Pertemuan}

\begin{longtable}{|c|>{\raggedright\arraybackslash}p{5.5cm}|>{\raggedright\arraybackslash}p{5.5cm}|}
\hline
\textbf{Mg} & \textbf{Sub-CPMK} & \textbf{Materi / Bahan Kajian} \\
\hline
\endfirsthead

\hline
\textbf{Mg} & \textbf{Sub-CPMK} & \textbf{Materi / Bahan Kajian} \\
\hline
\endhead

1 & Sub-CPMK 1 & Pengantar kriptografi, model keamanan, ancaman \\
\hline
2 & Sub-CPMK 2 & Aritmetika modular, operasi dasar, invers modular \\
\hline
3 & Sub-CPMK 2 & Algoritma Euclid dan Euclid diperluas (implementasi C) \\
\hline
4 & Sub-CPMK 2 & Teori bilangan: prima, fast exponentiation (implementasi C) \\
\hline
5 & Sub-CPMK 3 & Cipher Caesar dan substitusi sederhana (implementasi C) \\
\hline
6 & Sub-CPMK 3 & Cipher Vigenère dan polyalphabetic (implementasi C) \\
\hline
7 & Sub-CPMK 4 & Block cipher: DES, AES, mode ECB/CBC \\
\hline
8 & -- & \textbf{Ujian Tengah Semester (UTS)} \\
\hline
9 & Sub-CPMK 4 & Algoritma RSA, enkripsi asimetris (implementasi C) \\
\hline
10 & Sub-CPMK 5 & Fungsi hash (SHA, MD5), integritas data \\
\hline
11 & Sub-CPMK 5 & Tanda tangan digital (implementasi berbasis RSA) \\
\hline
12 & Sub-CPMK 6 & Protokol kriptografi, Diffie-Hellman, TLS \\
\hline
13 & Sub-CPMK 6 & Manajemen kunci, siklus hidup, distribusi \\
\hline
14 & Sub-CPMK 7 & Kriptanalisis: frequency analysis, Kasiski (praktikum) \\
\hline
15 & Sub-CPMK 8 & Presentasi proyek implementasi C \\
\hline
16 & -- & \textbf{Ujian Akhir Semester (UAS)} \\
\hline
\end{longtable}

\textit{Catatan: Mg = Minggu ke.}

\vspace{0.5cm}

% ============================================================
% 6. METODE PEMBELAJARAN
% ============================================================
\section*{6. Metode Pembelajaran}

Strategi atau pendekatan pembelajaran yang dipilih sesuai OBE yang menekankan aktivitas mahasiswa:

\begin{itemize}[leftmargin=*]
  \item \textbf{Ceramah Interaktif:} Penjelasan konsep matematika kriptografi dengan diskusi tanya jawab
  \item \textbf{Problem-Based Learning (PBL):} Mahasiswa menyelesaikan permasalahan keamanan informasi untuk mengidentifikasi kebutuhan kriptografi
  \item \textbf{Project-Based Learning:} Proyek implementasi algoritma kriptografi dalam bahasa C dengan dokumentasi dan pengujian
  \item \textbf{Praktikum Terbimbing:} Latihan coding C di laboratorium dengan bimbingan asisten
  \item \textbf{Diskusi Kasus:} Analisis studi kasus keamanan data dan protokol kriptografi
  \item \textbf{Flipped Classroom:} Mahasiswa mempelajari materi sebelum kelas, diskusi mendalam di kelas
  \item \textbf{Studi Kasus:} Analisis implementasi kriptografi pada aplikasi nyata (e-commerce, komunikasi aman)
\end{itemize}

\vspace{0.5cm}

% ============================================================
% 7. PENGALAMAN BELAJAR MAHASISWA
% ============================================================
\section*{7. Pengalaman Belajar Mahasiswa}

Deskripsi tugas, aktivitas, atau pengalaman belajar yang mendukung pencapaian Sub-CPMK:

\begin{itemize}[leftmargin=*]
  \item Mengimplementasikan cipher Caesar dan Vigenère dalam bahasa C
  \item Mengimplementasikan algoritma Euclid diperluas dan fast modular exponentiation dalam C
  \item Mengimplementasikan algoritma RSA dalam C (dengan bilangan kecil atau library seperti OpenSSL)
  \item Melakukan praktikum frequency analysis pada cipher substitusi
  \item Melakukan Kasiski examination untuk menyerang cipher Vigenère
  \item Merancang sistem enkripsi hibrida (AES + RSA) dalam C
  \item Membuat dokumentasi teknis dan laporan praktikum untuk setiap implementasi
  \item Berkolaborasi dalam tim untuk menyelesaikan proyek implementasi kriptografi
  \item Mempresentasikan hasil proyek dan analisis algoritma
  \item Melakukan pengujian kode dengan berbagai test cases
\end{itemize}

\vspace{0.5cm}

% ============================================================
% 8. KRITERIA, INDIKATOR, DAN BOBOT PENILAIAN
% ============================================================
\section*{8. Kriteria, Indikator, dan Bobot Penilaian}

Teknik/alat asesmen dipetakan ke Sub-CPMK/CPMK dengan bobot yang jelas:

\begin{longtable}{|>{\raggedright\arraybackslash}p{2.5cm}|>{\raggedright\arraybackslash}p{4cm}|>{\raggedright\arraybackslash}p{5.5cm}|c|}
\hline
\textbf{Komponen} & \textbf{Teknik Asesmen} & \textbf{Indikator/CPMK} & \textbf{Bobot (\%)} \\
\hline
\endfirsthead
\hline
\textbf{Komponen} & \textbf{Teknik Asesmen} & \textbf{Indikator/CPMK} & \textbf{Bobot (\%)} \\
\hline
\endhead
\hline
\endfoot

Tugas Individu & Coding Assignment (C) & Sub-CPMK 3.1, 3.2, 4.2 & 15 \\
\hline
Kuis & Multiple Choice \& Essay & Sub-CPMK 1.1, 1.2, 2.1, 2.2 & 10 \\
\hline
Praktikum & Lab Exercise (implementasi C) & Sub-CPMK 2.1, 2.2, 3.1, 7.1, 8.1 & 20 \\
\hline
UTS & Written Exam & CPMK-1, CPMK-2, CPMK-3 & 25 \\
\hline
Proyek & Implementasi C + Dokumentasi & CPMK-4, CPMK-5, CPMK-8 & 15 \\
\hline
UAS & Comprehensive Exam & Semua CPMK & 15 \\
\hline
\textbf{Total} & & & \textbf{100} \\
\hline
\end{longtable}

\textbf{Kriteria Penilaian:}
\begin{itemize}[leftmargin=*]
  \item A (85-100): Menguasai semua CPMK dengan sangat baik, mampu mengimplementasikan algoritma kriptografi dalam C
  \item B (70-84): Menguasai sebagian besar CPMK dengan baik
  \item C (60-69): Menguasai CPMK dasar dengan cukup
  \item D (50-59): Menguasai sebagian kecil CPMK
  \item E (<50): Belum menguasai CPMK yang ditetapkan
\end{itemize}

\vspace{0.5cm}

% ============================================================
% 9. EVALUASI DAN REFLEKSI PEMBELAJARAN (OPSIONAL)
% ============================================================
\section*{9. Evaluasi dan Refleksi Pembelajaran}

Penilaian sumatif/formatif untuk memantau ketercapaian outcome secara menyeluruh:

\begin{itemize}[leftmargin=*]
  \item \textbf{Evaluasi Formatif:} Kuis mingguan, latihan coding C, peer review untuk memberikan feedback berkelanjutan
  \item \textbf{Evaluasi Sumatif:} UTS dan UAS untuk mengukur pencapaian CPMK secara komprehensif
  \item \textbf{Refleksi Mahasiswa:} Jurnal belajar mingguan untuk refleksi terhadap pemahaman materi dan implementasi
  \item \textbf{Evaluasi Dosen:} Survey kepuasan mahasiswa di tengah dan akhir semester
  \item \textbf{Continuous Improvement:} Analisis hasil penilaian untuk perbaikan RPS di semester berikutnya
\end{itemize}

\vspace{0.5cm}

% ============================================================
% 10. DAFTAR REFERENSI
% ============================================================
\section*{10. Daftar Referensi}

Sumber belajar utama yang digunakan dalam penyusunan materi dan asesmen:

\begin{enumerate}[leftmargin=*]
  \item Stinson, D. R., \& Paterson, M. (2018). \textit{Cryptography: Theory and Practice} (4th ed.). CRC Press.
  
  \item Katz, J., \& Lindell, Y. (2020). \textit{Introduction to Modern Cryptography} (3rd ed.). CRC Press.
  
  \item Trappe, W., \& Washington, L. C. (2006). \textit{Introduction to Cryptography with Coding Theory} (2nd ed.). Pearson.
  
  \item Schneier, B. (1996). \textit{Applied Cryptography} (2nd ed.). John Wiley \& Sons.
  
  \item Stallings, W. (2017). \textit{Cryptography and Network Security: Principles and Practice} (7th ed.). Pearson.
  
  \item Menezes, A. J., van Oorschot, P. C., \& Vanstone, S. A. (1996). \textit{Handbook of Applied Cryptography}. CRC Press.
  
  \item Ferguson, N., Schneier, B., \& Kohno, T. (2010). \textit{Cryptography Engineering}. Wiley.
  
  \item NIST Cryptographic Standards. \url{https://csrc.nist.gov/}
  
  \item Cryptography Stack Exchange. \url{https://crypto.stackexchange.com/}
\end{enumerate}

\vspace{1cm}

\begin{flushright}
\begin{tabular}{c}
Disusun oleh,\\[2cm]
\textbf{[Nama Dosen]}\\
NIP. [NIP]
\end{tabular}
\end{flushright}

\end{document}
